% !TEX root = ../../../main.tex
\section{Systems of linear equations}
\label{sec:linear-algebra-i:systems-of-linear-equations}

% Sources: supplied lecture photos and September 17 handout, slides 3--17.
\subsection{Matrix form}

A system of \(n\) equations in \(m\) unknowns has the form
\[
  \sum_{j=1}^m a^i_jx^j=b^i,\qquad 1\leq i\leq n.
\]
Writing \(A=(a^i_j)\), \(x=(x^1,\ldots,x^m)^T\), and
\(b=(b^1,\ldots,b^n)^T\) gives
\[
  Ax=b,\qquad
  A\in\mathbb F^{n\times m},\quad x\in\mathbb F^m,\quad b\in\mathbb F^n.
\]
The number of equations and the number of unknowns need not agree.

\begin{example}
  The system
  \[
    2x-3y=5,\qquad x+y=-1,\qquad 3x+5y=0
  \]
  can be written as
  \[
    \begin{bmatrix}2&-3\\1&1\\3&5\end{bmatrix}
    \begin{bmatrix}x\\y\end{bmatrix}
    =\begin{bmatrix}5\\-1\\0\end{bmatrix}.
  \]
  Similarly, \(x+y+z=1\) and \(3y-5z=2\) become
  \[
    \begin{bmatrix}1&1&1\\0&3&-5\end{bmatrix}
    \begin{bmatrix}x\\y\\z\end{bmatrix}
    =\begin{bmatrix}1\\2\end{bmatrix}.
  \]
\end{example}

Over \(\R\), the equation \(x=b\) has exactly one solution, whereas
\[
  \begin{bmatrix}1&0\end{bmatrix}
  \begin{bmatrix}x\\y\end{bmatrix}=b
\]
has infinitely many, since \(y\) is arbitrary. The equations \(x=1\) and
\(x=0\) are inconsistent:
\[
  \begin{bmatrix}1\\1\end{bmatrix}x=\begin{bmatrix}1\\0\end{bmatrix}.
\]
Two equations in two unknowns may have no solution, one solution, or
infinitely many solutions. For example, \(I_2(x,y)^T=(a,b)^T\) has the
unique solution \(x=a\), \(y=b\).
% The earlier photo's claim that solvability requires equal equation and
% unknown counts contradicts these supplied rectangular examples.

\subsection{Augmented matrices and row operations}

Appending the right-hand side produces the augmented matrix
\([A\mid b]\). The equation \(Ax=b\) is equivalent to
\[
  [A\mid b]\begin{bmatrix}x\\-1\end{bmatrix}=0.
\]
The final coordinate here is fixed at \(-1\).

The elementary row operations exchange two rows, multiply a row by a
nonzero scalar, or add a multiple of one row to another. Apply an operation
to the entire augmented row, including its right-hand side. Reordering
variables corresponds to reordering the columns of \(A\).
% The handout says "rescale"; nonzero is made explicit to distinguish a
% reversible row operation from erasing an equation.

\begin{example}
  Consider
  \[
    2x-4y+6z=12,\qquad 3x+12y+24z=36.
  \]
  Dividing the first equation by \(2\) gives \(x-2y+3z=6\).
  Subtracting three times this equation from the second gives
  \(18y+15z=18\), or \(y+\frac56z=1\). Adding twice this last equation
  to the first eliminates \(y\), leaving
  \[
    x+\frac{14}{3}z=8,\qquad y+\frac56z=1.
  \]
  Thus all solutions are
  \[
    x=8-\frac{14}{3}z,\qquad y=1-\frac56z,\qquad z\in\R.
  \]
\end{example}

\begin{example}
  The system
  \[
    \begin{aligned}
      x^1+2x^2+3x^3+5x^4+x^5&=1,\\
      x^3+5x^4&=2,\\
      -3x^5&=9
    \end{aligned}
  \]
  has augmented matrix
  \[
    \left[\begin{array}{rrrrr|r}
      1&2&3&5&1&1\\
      0&0&1&5&0&2\\
      0&0&0&0&-3&9
    \end{array}\right].
  \]
  Divide row \(3\) by \(-3\), subtract the new row \(3\) from row \(1\),
  and then subtract three times row \(2\) from row \(1\). The successive
  matrices are
  \[
    \left[\begin{array}{rrrrr|r}
      1&2&3&5&1&1\\0&0&1&5&0&2\\0&0&0&0&1&-3
    \end{array}\right],
    \qquad
    \left[\begin{array}{rrrrr|r}
      1&2&3&5&0&4\\0&0&1&5&0&2\\0&0&0&0&1&-3
    \end{array}\right],
  \]
  and finally
  \[
    \left[\begin{array}{rrrrr|r}
      1&2&0&-10&0&-2\\0&0&1&5&0&2\\0&0&0&0&1&-3
    \end{array}\right].
  \]
\end{example}

\subsection{Triangular systems and factorization}

\begin{example}
  An upper triangular system can be solved from the last equation upward:
  \[
    2x+y+z=3,\qquad y-3z=4,\qquad -z=2.
  \]
  In augmented form, divide the last row by \(-1\), then eliminate \(z\)
  from the first two rows:
  \[
    \left[\begin{array}{rrr|r}
      2&1&1&3\\0&1&-3&4\\0&0&-1&2
    \end{array}\right]
    \longrightarrow
    \left[\begin{array}{rrr|r}
      2&1&1&3\\0&1&-3&4\\0&0&1&-2
    \end{array}\right]
    \longrightarrow
    \left[\begin{array}{rrr|r}
      2&1&0&5\\0&1&0&-2\\0&0&1&-2
    \end{array}\right].
  \]
  Subtract row \(2\) from row \(1\), then divide row \(1\) by \(2\):
  \[
    \left[\begin{array}{rrr|r}
      2&0&0&7\\0&1&0&-2\\0&0&1&-2
    \end{array}\right]
    \longrightarrow
    \left[\begin{array}{rrr|r}
      1&0&0&7/2\\0&1&0&-2\\0&0&1&-2
    \end{array}\right].
  \]
  The solution is \(x=\frac72\), \(y=-2\), \(z=-2\).
\end{example}

For a square system \(Ax=b\), an available factorization \(A=LU\), with
\(L\) lower triangular and \(U\) upper triangular, reduces the problem to
\[
  Ly=b,\qquad Ux=y.
\]
Indeed, \(Ax=LUx=Ly=b\). Sparse matrices have many zero entries.
Their pattern of zeros can support more efficient solution strategies.
Fast matrix multiplication is another topic in numerical linear algebra.

\subsection{Normalized systems}

After row operations and a possible reordering of variables, consider the
following four forms.

\begin{enumerate}
  \item For \(I_mx=b'\), every variable is determined:
    \(x^j=(b')^j\), \(1\leq j\leq m\).
  \item For
    \[
      [I_k\mid C]x=b',\qquad C\in\mathbb F^{k\times(m-k)},\quad k<m,
    \]
    the last \(m-k\) variables are free. The first \(k\) satisfy
    \[
      x^j+\sum_{\ell=k+1}^m c^j_\ell x^\ell=(b')^j,\qquad 1\leq j\leq k.
    \]
  \item For
    \[
      \begin{bmatrix}I_k\\0_{(n-k)\times k}\end{bmatrix}x=b',
    \]
    the equations require \(x^j=(b')^j\) for \(1\leq j\leq k\) and
    \((b')^{k+1}=\cdots=(b')^n=0\). A nonzero entry among these last
    right-hand sides makes the system inconsistent.
  \item In the combined case,
    \[
      \begin{bmatrix}
        I_k&C\\
        0_{(n-k)\times k}&0_{(n-k)\times(m-k)}
      \end{bmatrix}x=b',
    \]
    there are free variables \(x^{k+1},\ldots,x^m\), while the last
    \(n-k\) equations require \((b')^{k+1}=\cdots=(b')^n=0\).
\end{enumerate}
% September 17 slides 16--17 resolve the ambiguous indices in the earlier
% photographed third normalization case. C denotes the remaining block,
% rather than reusing A for both the original matrix and its reduced block.
